% Entwurf: Kapitel 1 (Einleitung/Motivation) und Kapitel 3 (Anforderungsanalyse)
% Bachelorarbeit Sergi Koniashvili, Universität Bremen
% Stand: 2026-11-01 — Diskussionsgrundlage für die Besprechung am 2026-11-02.
% Fakten ausschließlich aus docs/thesis/THESIS_KNOWLEDGE_BASE.md (VERIFIED).
% Nicht belegte Stellen sind mit [TODO] markiert.
\documentclass[11pt,a4paper]{article}
\usepackage{fontspec}
\usepackage[ngerman]{babel}
\usepackage[a4paper,margin=2.5cm]{geometry}
\usepackage{csquotes}
\usepackage{booktabs}
\usepackage{longtable}
\usepackage{array}
\usepackage{xcolor}
\usepackage{enumitem}
\usepackage[hidelinks]{hyperref}

\newcommand{\todo}[1]{\textcolor{red!70!black}{[TODO: #1]}}
\newcommand{\frage}[1]{\par\smallskip\noindent\colorbox{blue!8}{\parbox{\dimexpr\linewidth-2\fboxsep}{\small\textbf{Frage an die Betreuer:} #1}}\par\smallskip}

\title{Entwurf: Motivation und literaturbasierte Anforderungsanalyse\\[0.3em]
\large LLM-gestützte Integration von Klima- und Gesundheitsdaten}
\author{Sergi Koniashvili\\ Bachelorarbeit Informatik, Universität Bremen}
\date{Stand: 1.~November 2026 -- Diskussionsgrundlage}

\begin{document}
\maketitle

\noindent\textit{Hinweis:} Dieser Entwurf zeigt den geplanten Aufbau sowie erste
Fassungen von Kapitel~1 und Kapitel~3. Rot markierte Stellen sind noch nicht
belegt bzw.\ offen. Zitierstil vorläufig Autor-Jahr (siehe Frage in
Abschnitt~\ref{sec:fragen}).

\section*{Geplante Gliederung der Arbeit}
\begin{enumerate}[nosep]
  \item Einleitung (Motivation, Problemstellung, Zielsetzung, Aufbau)
  \item Grundlagen (Klima und vektorübertragene Krankheiten; Datenquellen:
        Fallzahlen, Reanalyse-Klimadaten, Bevölkerung; ETL; LLMs)
  \item Anforderungsanalyse (Ableitung aus Beispielpublikationen)
  \item Verwandte Arbeiten und Werkzeuge \todo{Werkzeuge aus offizieller Dokumentation prüfen}
  \item Konzept und Architektur
  \item Implementierung
  \item Evaluation (Fallstudien CS1--CS3)
  \item Diskussion, Grenzen und Ausblick
  \item Fazit
\end{enumerate}

% ---------------------------------------------------------------------------
\section{Einleitung}

\subsection{Motivation}

Die Wechselwirkung zwischen Klima und vektorübertragenen Krankheiten wie Dengue
wird zunehmend mit datenbasierten Modellen untersucht. Ein aktuelles Beispiel
ist die Studie von Suttirat et al.\ (2026), die monatliche Fallzahlen des
Dengue-hämorrhagischen Fiebers (DHF) für 77 Provinzen Thailands über 20 Jahre
(2003--2022) mit Klima-, Fernerkundungs- und sozioökonomischen Daten verknüpft
und mit XGBoost und SHAP analysiert (S.~5--8). Lowe et al.\ (2021) untersuchen
auf ähnliche Weise den Einfluss hydrometeorologischer Extreme und der
Urbanisierung auf das Dengue-Risiko in 558 Mikroregionen Brasiliens
(S.~e211--e212).

Beiden Studien ist gemeinsam, dass vor der eigentlichen Modellierung eine
umfangreiche Datenvorbereitung steht. Suttirat et al.\ beschreiben diese
Schritte wie folgt:

\begin{quote}
\enquote{Gridded climatic and remote-sensing data were spatially aggregated to
the provincial level; lagged variables and bioclimatic indices were then derived
from monthly provincial temperature and precipitation series. Socioeconomic data
were obtained directly at the provincial level: the biennially reported household
income series was linearly interpolated to fill missing years, and all
socioeconomic series were then upsampled to monthly resolution by assigning each
year's reported value (typically released at year-end) to every month of that
year.} (Suttirat et al., 2026, S.~4)
\end{quote}

Diese Schritte sind nicht spezifisch für eine einzelne Studie. Lowe et al.\
(2021) führen weitgehend dieselben Schritte durch: Sie beziehen monatliche
Fallzahlen aus dem nationalen Meldesystem, berechnen mit jährlichen
Bevölkerungsschätzungen die Inzidenz pro 100\,000 Einwohner, aggregieren
gerasterte Klimadaten (Tmin, Tmax, PDSI) flächengewichtet auf die
Untersuchungseinheiten und laden Klimadaten für ein zusätzliches Jahr vor der
ersten Fallbeobachtung, \enquote{to allow for a lag period before the first
dengue observation in January, 2001} (S.~e211). Auch Suttirat et al.\ erzeugen
zeitverzögerte Variablen (Lags von 1--3 Monaten) für alle Klima- und
Fernerkundungsvariablen (S.~7). Tabelle~\ref{tab:schritte} stellt die Schritte
beider Studien gegenüber.

Die Daten selbst sind oft öffentlich verfügbar, aber nicht unmittelbar
analysebereit. Lowe et al.\ nutzen ausschließlich offene Daten (S.~e218), jedoch
aus fünf heterogenen Quellen mit unterschiedlichen Formaten und Auflösungen
(monatlich, jährlich, statisch, 0,5°-Raster); die Zusammenstellung der Datenbank
wird in den Autorenbeiträgen als eigene Aufgabe genannt (S.~e218). Für die
Gesundheitsdaten halten Clarke et al.\ (2024) bei der Beschreibung des
OpenDengue-Datensatzes fest:

\begin{quote}
\enquote{Dengue case data exists in multiple formats from a wide variety of
sources that require various processing methodologies. [\ldots] Locating,
extracting, processing and standardising data all takes time but is essential to
enable reuse and re-analysis.} (Clarke et al., 2024, S.~2)
\end{quote}

OpenDengue löst damit die Sammlung der Fallzahlen. Die Verknüpfung mit
Klimadaten, die Clarke et al.\ als einen Anwendungsfall nennen (S.~1, 9),
bleibt jedoch Aufgabe der Forschenden.

Schließlich setzt dieser Arbeitsablauf Programmierkenntnisse voraus: Suttirat
et al.\ führten die Datenaufbereitung mit Python (pandas) durch und berechneten
bioklimatische Indikatoren mit dem R-Paket \texttt{dismo} (S.~5, 8). Für
Forschende ohne Programmierkenntnisse, etwa im Bereich Public Health, stellt
dies eine Hürde dar \todo{Quelle zu fehlenden Programmier-/Data-Science-Kenntnissen im Public-Health-Bereich}.
Bestehende Werkzeuge decken diese Schritte nur teilweise ab
\todo{Verwandte Arbeiten: Copernicus CDS, Google Earth Engine, DHIS2 Climate \& Health, IRI ENACTS -- jeweils aus offizieller Dokumentation belegen}.

\subsection{Problemstellung und Zielsetzung}

Ziel dieser Arbeit ist ein System, das Forschenden die der Modellierung
vorgelagerte Datenvorbereitung abnimmt -- Quellensuche, Bereinigung,
Harmonisierung, Zusammenführung und Dokumentation --, und zwar ohne
Programmierkenntnisse und nachvollziehbar, sodass sie sich auf das eigentliche
Modell konzentrieren können. Die Forschenden beschreiben den benötigten
Datensatz (Krankheit, Region, Zeitraum, Klimavariablen) in einem Formular oder
in natürlicher Sprache; ein Large Language Model (LLM) überführt die Anfrage in
eine strukturierte Integrationsanfrage. Das System wählt geeignete Quellen,
führt die Daten zu einer analysebereiten Tabelle zusammen, berechnet abgeleitete
Größen (Inzidenz, zeitverzögerte Variablen) und protokolliert jeden Schritt.

Das System ist ausdrücklich \emph{kein} Vorhersagemodell: Es bereitet
Eingangsdaten für nachgelagerte datenbasierte Modelle vor. Ob es den manuellen
Aufwand tatsächlich reduziert, wird in der Evaluation gemessen
(Kapitel~7, Fallstudie CS1: derselbe Zieldatensatz wird manuell mit
Python/pandas und mit dem System erstellt; erfasst werden Zeit, Anzahl manueller
Schritte und Übereinstimmung der Ergebnisse).

\subsection{Aufbau der Arbeit}
\todo{nach Festlegung der Gliederung ausformulieren}

% ---------------------------------------------------------------------------
\setcounter{section}{2}
\section{Anforderungsanalyse}

\subsection{Vorgehen}

Die Anforderungen werden aus Beispielpublikationen der datenbasierten
Klima-Gesundheits-Modellierung abgeleitet. Für jede Publikation werden die vor
der Modellierung durchgeführten Datenvorbereitungsschritte aus dem Methodenteil
extrahiert. Jeder Schritt, der für die Bereitstellung eines analysebereiten
Datensatzes notwendig ist, wird in eine Anforderung überführt oder begründet
ausgeschlossen (Abgrenzung). Die Zuordnung Schritt $\rightarrow$ Publikation
(mit Seitenangabe) $\rightarrow$ Anforderung wird in einer Traceability-Matrix
festgehalten (Tabelle~\ref{tab:trace}).

Grundlage sind derzeit drei Publikationen:
\begin{itemize}[nosep]
  \item \textbf{Suttirat et al.\ (2026)}: Dengue, Thailand, 77 Provinzen,
        2003--2022, XGBoost/SHAP -- Hauptfallstudie;
  \item \textbf{Lowe et al.\ (2021)}: Dengue, Brasilien, 558 Mikroregionen,
        2001--2019, bayesianisches hierarchisches Modell -- zeigt, dass die
        Schritte typisch sind;
  \item \textbf{Clarke et al.\ (2024)}: OpenDengue-Datensatz (Data Descriptor) --
        Anforderungen an Herkunft, Version und Vergleichbarkeit der Fallzahlen.
\end{itemize}

\subsection{Datenvorbereitungsschritte in den Beispielpublikationen}

\begin{longtable}{@{}p{0.27\linewidth}p{0.33\linewidth}p{0.33\linewidth}@{}}
\caption{Datenvorbereitung vor der Modellierung (Auszug)}\label{tab:schritte}\\
\toprule
Schritt & Suttirat et al.\ (2026) & Lowe et al.\ (2021) \\
\midrule
\endfirsthead
\toprule
Schritt & Suttirat et al.\ (2026) & Lowe et al.\ (2021) \\
\midrule
\endhead
Fallzahlen beziehen & monatliche DHF-Fälle, Bureau of Epidemiology (S.~5) &
  monatliche Meldefälle, SINAN/DATASUS (S.~e211) \\
Verwaltungseinheiten harmonisieren & Bueng Kan in Nong Khai zusammengeführt (S.~5) & -- \\
Klimavariablen beziehen & TerraClimate: u.\,a.\ Niederschlag, Tmin, Tmax (S.~5) &
  CRU TS 4.04: Tmin, Tmax, PDSI (S.~e211) \\
Räumliche Aggregation & Raster auf Provinzebene gemittelt, Methode nicht angegeben (S.~5) &
  flächengewichtetes Mittel mit \texttt{exactextractr} (S.~e211) \\
Zeitverzögerte Variablen & Lags 1--3 Monate (S.~7) & Lags 0--6 Monate (DLNM, S.~e212) \\
Puffer für Lags & -- & Klimadaten ein Jahr vor Beginn der Fallreihe (S.~e211) \\
Inzidenz pro 100\,000 & ja (S.~7) & ja, mit jährlicher Bevölkerung (S.~e211) \\
Jahreswerte $\rightarrow$ Monate & Jahreswert jedem Monat zugeordnet (S.~4--5) &
  nicht beschrieben \\
Fehlende Jahre auffüllen & Haushaltseinkommen linear interpoliert (S.~5) & -- \\
Abgeleitete Indizes & 19 bioklimatische Indikatoren (\texttt{dismo}, S.~5) & -- \\
Werkzeuge & Python/pandas, R (S.~5, 8) & R, \texttt{exactextractr}, R-INLA (S.~e211--e212) \\
\bottomrule
\end{longtable}

\subsection{Abgeleitete Anforderungen (Traceability-Matrix)}

\begin{longtable}{@{}p{0.27\linewidth}p{0.22\linewidth}p{0.33\linewidth}p{0.1\linewidth}@{}}
\caption{Zuordnung Schritt $\rightarrow$ Publikation $\rightarrow$ Anforderung}\label{tab:trace}\\
\toprule
Schritt / Beobachtung & Beleg & Anforderung & Status \\
\midrule
\endfirsthead
\toprule
Schritt / Beobachtung & Beleg & Anforderung & Status \\
\midrule
\endhead
Inzidenz pro 100\,000 als Zielgröße & Suttirat S.~7; Lowe S.~e211 &
  FR-17: Bevölkerungsdaten anbinden, jährliche Werte auf Monate interpolieren, Inzidenz berechnen & umgesetzt \\
Zeitverzögerte Klimavariablen & Suttirat S.~7 (1--3); Lowe S.~e212 (0--6) &
  FR-23: Lag-Spalten 1--6 im Datensatz & umgesetzt \\
Klimadaten vor Beginn der Fallreihe & Lowe S.~e211 &
  FR-23: Lag-Puffer, damit keine Anfangszeilen verloren gehen & umgesetzt \\
Tmin/Tmax statt nur Mitteltemperatur & Suttirat S.~5; Lowe S.~e211 &
  FR-24: Tmin/Tmax als Mittel der täglichen Minima/Maxima & umgesetzt \\
Unterschiedliche zeitliche Auflösungen & Suttirat S.~4--5 &
  FR-20: Auflösungsharmonisierung \todo{Unterschied diskutieren: Suttirat wiederholt Jahreswerte (Bestandsgröße), FR-20 erhält Summen (Zählgröße)} & umgesetzt \\
Herkunft muss nachvollziehbar sein; Version angeben & Clarke S.~2, 4, 12 &
  FR-25: Export mit Metadaten (Quellen, Lizenzen, Version, Schritte) & umgesetzt \\
Räumliche Abdeckung variiert zwischen Ländern (subnational nur für 40 von 102 Ländern) & Clarke S.~1, 8 &
  FR-26: Datenlücken erkennen und begründen, nicht automatisch füllen & umgesetzt \\
Vergleichbarkeit zwischen Quellen beurteilen & Clarke S.~2 &
  FR-27: zwei Quellen im Überlappungszeitraum vergleichen & umgesetzt \\
Zusammenhänge explorativ prüfen & Lags: Suttirat S.~7; Lowe S.~e212 &
  FR-19: Kreuzkorrelationen Lag 0--3 (Pearson, Spearman) & umgesetzt \\
\bottomrule
\end{longtable}

\noindent Zu FR-19 ist eine Grenze festzuhalten: Korrelationen erfassen keine
nichtlinearen Schwelleneffekte, wie sie Suttirat et al.\ fanden (etwa
21\,°C für Tmin mit Lag~1 und 32\,°C für Tmax mit Lag~3, S.~12); Lowe et al.\
betonen, dass der Einfluss des Klimas \enquote{complex, non-linear, and often
delayed} ist (S.~e211). Die Korrelationen dienen daher nur der explorativen
Datenprüfung, nicht der Modellierung.

\subsection{Anforderungen ohne direkten Literaturbezug}

Der ursprüngliche Anforderungskatalog (FR-01 bis FR-22) entstand vor der
Literaturanalyse. Ein Teil davon lässt sich nachträglich begründen (FR-17,
FR-19, FR-20, siehe oben). Weitere Anforderungen ergeben sich aus dem Ziel
\enquote{ohne Programmierkenntnisse} (natürlichsprachliche Anfrage FR-01,
verständliche Erklärungen FR-09, Chat-Assistent FR-13, Weltkarte FR-22,
PDF-Bericht FR-21) und aus Vorgaben der Betreuer (u.\,a.\ FR-03 bis FR-08).
\frage{Sollen diese Anforderungen als Usability-Anforderungen mit dem Ziel
\enquote{ohne Programmierung} begründet oder aus der Arbeit gestrichen werden?}

\subsection{Abgrenzung}

Folgende Schritte aus den Beispielpublikationen werden bewusst nicht umgesetzt;
sie werden in Kapitel~8 als Grenzen diskutiert:
\begin{itemize}[nosep]
  \item subnationale Einheiten und räumliche Aggregation von Rasterdaten
        (Suttirat S.~5; Lowe S.~e211) -- das System nutzt einen repräsentativen
        Punkt pro Land;
  \item Harmonisierung sich ändernder Verwaltungseinheiten (Suttirat S.~5);
  \item bioklimatische Indikatoren und Fernerkundungsdaten (Suttirat S.~5);
  \item automatisches Auffüllen von Lücken aus einer anderen Quelle: Fallzahlen
        verschiedener Quellen beruhen auf unterschiedlichen Falldefinitionen --
        OpenDengue zählt alle gemeldeten Fälle (Clarke S.~5), Suttirat et al.\
        nur DHF (S.~5) --, sodass eine verdeckte Mischung die Vergleichbarkeit
        beeinträchtigen würde. Das System erkennt Lücken (FR-26) und bietet einen
        expliziten Quellenvergleich an (FR-27).
\end{itemize}

% ---------------------------------------------------------------------------
\section*{Offene Fragen für die Besprechung}\label{sec:fragen}
\begin{enumerate}[nosep]
  \item Entspricht die Argumentationskette der Motivation (Modelle $\rightarrow$
        aufwendige, typische Datenvorbereitung $\rightarrow$ Programmierkenntnisse
        nötig $\rightarrow$ System) Ihrer Vorstellung?
  \item Reichen zwei Modellierungsstudien und ein Data Descriptor, oder sollen
        weitere Studien (z.\,B.\ Malaria) aufgenommen werden?
  \item Ist die Traceability-Matrix als Form der Anforderungsableitung geeignet?
  \item Welcher Zitierstil wird erwartet (Autor-Jahr oder numerisch)?
  \item Gliederung: Passt die Position von \enquote{Verwandte Arbeiten} nach der
        Anforderungsanalyse?
\end{enumerate}

\section*{Literatur (Entwurf)}
\begin{small}
\begin{description}[nosep,leftmargin=1em,labelindent=0pt]
\item Clarke, J., Lim, A., Gupte, P., Pigott, D.\,M., van Panhuis, W.\,G., \&
  Brady, O.\,J. (2024). A global dataset of publicly available dengue case count
  data. \emph{Scientific Data}, 11, 296.
  \url{https://doi.org/10.1038/s41597-024-03120-7}
\item Lowe, R., Lee, S.\,A., O'Reilly, K.\,M., Brady, O.\,J., Bastos, L.,
  Carrasco-Escobar, G., et al.\ (2021). Combined effects of hydrometeorological
  hazards and urbanisation on dengue risk in Brazil: a spatiotemporal modelling
  study. \emph{The Lancet Planetary Health}, 5(4), e209--e219.
  \url{https://doi.org/10.1016/S2542-5196(20)30292-8}
\item Suttirat, P., Chadsuthi, S., Aekthong, S., Rocklöv, J., Bicout, D.\,J.,
  Haddawy, P., Su Yin, M., Lawpoolsri, S., \& Modchang, C. (2026). Machine
  learning reveals temperature as a key predictor of dengue risk across
  Thailand's provinces: A 20-year analysis. \emph{PLoS Neglected Tropical
  Diseases}, 20(7), e0014590.
  \url{https://doi.org/10.1371/journal.pntd.0014590}
\end{description}
\end{small}

\end{document}
